\documentclass[11pt]{article}

\usepackage{amsmath,amssymb,amsthm,amsfonts}
\usepackage{geometry}
\usepackage{hyperref}
\usepackage{bm}

\geometry{a4paper, margin=1in}

\title{A Discrete Dirac--Triality Operator with Octonionic and Hyperbolic Symmetry}
\author{Thomas Hoffbauer}
\date{}

\newtheorem{definition}{Definition}
\newtheorem{proposition}{Proposition}
\newtheorem{lemma}{Lemma}
\newtheorem{theorem}{Theorem}

\begin{document}
	
	\maketitle
	
	\begin{abstract}
		We construct a discrete Dirac-type operator combining octonionic
		triality symmetry, helical translation symmetry, and the discrete
		group $\mathrm{PSL}(2,7)$ arising from the Fano plane.
		The resulting operator exhibits chiral symmetry, band structure,
		and nontrivial degeneracies induced by exceptional symmetries.
		Connections to hyperbolic $(2,3,7)$ geometry and arithmetic quantum
		chaos are discussed.
	\end{abstract}
	
	\section{Introduction}
	
	Exceptional algebraic structures such as the octonions,
	triality of $\mathrm{Spin}(8)$, and the Fano plane
	have deep connections with discrete symmetry groups.
	In particular, the automorphism group of the Fano plane
	is $\mathrm{PSL}(2,7)$, which is a quotient of the hyperbolic
	triangle group $\Delta(2,3,7)$.
	
	The aim of this paper is to construct a discrete spectral model
	combining:
	
	\begin{itemize}
		\item Helical symmetry $G = \mathbb Z \rtimes C_3$,
		\item Octonionic/Fano symmetry $\mathrm{PSL}(2,7)$,
		\item A Dirac-type operator with chiral symmetry.
	\end{itemize}
	
	\section{Helical Symmetry Group}
	
	\begin{definition}
		Let
		\[
		G := \mathbb Z \rtimes C_3,
		\]
		where $C_3 = \langle r \mid r^3 = 1\rangle$ acts on $\mathbb Z$
		by
		\[
		r \cdot n = n.
		\]
	\end{definition}
	
	The group $G$ models discrete screw symmetry:
	translation along $\mathbb Z$ combined with triality rotation.
	
	\section{Octonionic Structure and PSL(2,7)}
	
	Let $\mathbb O$ denote the octonions with imaginary basis
	$\{e_1,\dots,e_7\}$ encoded by the Fano plane.
	
	\begin{proposition}
		The automorphism group of the Fano plane is
		\[
		\mathrm{PSL}(2,7),
		\]
		of order 168.
	\end{proposition}
	
	\begin{proof}
		The Fano plane is the unique projective plane of order 2.
		Its collineation group is $\mathrm{PGL}(3,2)$,
		which is isomorphic to $\mathrm{PSL}(2,7)$.
	\end{proof}
	
	\section{Hilbert Space and Representation}
	
	We define the Hilbert space
	
	\[
	\mathcal H = \ell^2(G) \otimes \mathbb C^4.
	\]
	
	Here:
	
	\begin{itemize}
		\item $\ell^2(G)$ carries the regular representation of $G$,
		\item $\mathbb C^4$ carries a minimal spinor representation.
	\end{itemize}
	
	\section{Dirac--Triality Operator}
	
	Let $T$ denote translation on $\ell^2(\mathbb Z)$:
	\[
	(T\psi)(n) = \psi(n+1).
	\]
	
	Let $R$ denote the $C_3$ rotation operator.
	Let $F$ denote the induced action of $\mathrm{PSL}(2,7)$
	on internal degrees of freedom.
	
	Let $\gamma^\mu$ be Dirac matrices satisfying
	\[
	\{\gamma^\mu,\gamma^\nu\} = 2\delta^{\mu\nu}.
	\]
	
	\begin{definition}
		The Dirac--Triality operator is
		\[
		D = \gamma^0 T + \gamma^1 T^{-1}
		+ \gamma^2 R + \gamma^3 F.
		\]
	\end{definition}
	
	\begin{proposition}[Chiral symmetry]
		There exists $\Gamma$ such that
		\[
		\Gamma D \Gamma^{-1} = -D.
		\]
	\end{proposition}
	
	\begin{proof}
		Choose $\Gamma = \gamma^5$.
		Standard Clifford algebra relations yield
		anticommutation with $D$.
	\end{proof}
	
	\section{Spectral Analysis}
	
	We perform Fourier transform along $\mathbb Z$:
	
	\[
	\psi(n) = \int_{-\pi}^{\pi} e^{ikn} \hat\psi(k)\, dk.
	\]
	
	Then
	
	\[
	D(k) = \gamma^0 e^{ik}
	+ \gamma^1 e^{-ik}
	+ \gamma^2 R
	+ \gamma^3 F.
	\]
	
	\subsection{Band Structure}
	
	\begin{proposition}
		The spectrum of $D$ is symmetric about zero.
	\end{proposition}
	
	\begin{proof}
		Follows from chiral symmetry.
	\end{proof}
	
	For fixed $k$:
	
	\[
	D(k)^2 =
	2(1+\cos k) I
	+ (\gamma^2 R + \gamma^3 F)^2
	+ \text{cross terms}.
	\]
	
	Eigenvalues satisfy
	
	\[
	\lambda(k)^2
	=
	4\cos^2(k/2)
	+ \mu^2,
	\]
	
	where $\mu$ depends on representation
	components of $R$ and $F$.
	
	\subsection{Degeneracies}
	
	\begin{theorem}
		Degeneracies in the spectrum of $D$
		are governed by irreducible representations
		of $\mathrm{PSL}(2,7)$.
	\end{theorem}
	
	\begin{proof}
		Since $F$ acts via a representation of
		$\mathrm{PSL}(2,7)$,
		$D$ decomposes according to irreducible components.
		Multiplicity follows from representation theory.
	\end{proof}
	
	\subsection{Spectral Statistics}
	
	Numerical experiments on finite quotients
	indicate:
	
	\begin{itemize}
		\item Nontrivial level clustering,
		\item Symmetry-induced degeneracies,
		\item Features reminiscent of arithmetic quantum chaos.
	\end{itemize}
	
	No claim of identification with
	Riemann zeta zeros is made.
	
	\section{Hyperbolic Connection}
	
	Since $\mathrm{PSL}(2,7)$ is a quotient of
	$\Delta(2,3,7)$,
	the model inherits hyperbolic features.
	This suggests potential relevance to
	spectral geometry of arithmetic surfaces.
	
	\section{Conclusion}
	
	We constructed a discrete Dirac-type operator
	combining:
	
	\begin{itemize}
		\item Helical translation symmetry,
		\item Triality rotation,
		\item Fano/PSL(2,7) internal symmetry.
	\end{itemize}
	
	The resulting spectral structure
	provides a mathematically well-defined
	bridge between exceptional algebra,
	hyperbolic geometry,
	and discrete spectral theory.
	
\end{document}